% Generated draft: requires mathematical and source review.
\documentclass[11pt]{article}
\usepackage[margin=25mm]{geometry}
\usepackage{fontspec}
\setmainfont{Latin Modern Roman}
\usepackage{amsmath,amssymb,mathtools,unicode-math}
\setmathfont{Latin Modern Math}
\usepackage{enumitem,needspace,xurl,hyperref}
\setlength{\parindent}{0pt}
\setlength{\parskip}{5pt}
\setlength{\emergencystretch}{3em}
\newcommand{\N}{\mathbb{N}}
\newcommand{\Z}{\mathbb{Z}}
\newcommand{\Q}{\mathbb{Q}}
\newcommand{\R}{\mathbb{R}}
\newcommand{\C}{\mathbb{C}}
\newcommand{\F}{\mathbb{F}}
\newcommand{\A}{\mathbb{A}}
\newcommand{\PP}{\mathbb{P}}
\newcommand{\E}{\mathbb{E}}
\newcommand{\eps}{\varepsilon}
\newcommand{\dd}{\,\mathrm{d}}
\begin{document}
\section*{UnsolvedMath ID: 30003500}
\textbf{Source problem code:} catalogue entry
\subsection{Definitions and mathematical statement}
Let \(X\subseteq\R^d\) be closed. Interpret \emph{fat} in its usual topological sense:
\[
X=\overline{\operatorname{int}_{\R^d}(X)}.
\]
The supplied record asks for conditions on such sets under which
\[
A^\infty(X)=C^\infty(X)
\qquad\text{and}\qquad
A^M(X)=C^M(X).
\]
It does not define these four classes or specify the parameter \(M\).

Under the conventional real-valued, ambient-extension interpretation, the smooth classes can be defined by
\[
C^\infty(X)=\left\{f:X\to\R:\ \exists\,U\subseteq\R^d\text{ open},\ X\subseteq U,\ \exists\,F\in C^\infty(U),\ F|_X=f\right\}
\]
and
\[
A^\infty(X)=\left\{f:X\to\R:\ f\circ c\in C^\infty(\R)\text{ for every }c\in C^\infty(\R,\R^d)\text{ with }c(\R)\subseteq X\right\}.
\]
With these explicitly stated conventions, the smooth question is to characterize the closed fat sets \(X\) for which every function passing all these smooth-curve tests extends smoothly to an open neighborhood of \(X\). This is a question about conditions on \(X\), not an assertion that the equality holds for every closed fat set.

For the ultradifferentiable question, suppose that a parameter \(M\) and classes \(C^M(U)\) on open sets have first been specified. A possible corresponding interpretation is
\[
C^M(X)=\left\{f:X\to\R:\ \exists\,U\subseteq\R^d\text{ open},\ X\subseteq U,\ \exists\,F\in C^M(U),\ F|_X=f\right\},
\]
\[
A^M(X)=\left\{f:X\to\R:\ f\circ c\in C^M(\R)\text{ for every }c\in C^M(\R,\R^d)\text{ with }c(\R)\subseteq X\right\},
\]
where the coordinates of a \(C^M\) curve belong to \(C^M(\R)\). Under this interpretation, ask which closed fat sets have the indicated equality, for the specified \(M\). These definitions are conditional interpretations requiring confirmation: the record does not supply the weight data, growth bounds, Roumieu or Beurling convention, restrictions on \(M\), curve-testing convention, or intended extension notion. Consequently, it does not determine a unique precise ultradifferentiable problem.

The statement concerns closed fat sets generally. The background separately refers to a general compact-set question; compactness is not imposed in the displayed source statement.
\subsection{Short English statement}
Determine geometric conditions on a closed fat set under which smoothness, or a specified form of ultradifferentiability, can be detected by composing functions with curves in the set. The smooth formulation above uses an explicit conventional interpretation. The ultradifferentiable formulation remains conditional because the source does not define \(M\) or the relevant function classes.
\subsection{Sources}
\begin{itemize}
\item Supplied local UnsolvedMath registry record: permanent ID \(30003500\), title \emph{Arc-Smoothness on Closed Fat Sets}, exact source problem code \texttt{catalogue entry}, in the Oberwolfach Reports Open Problems collection. The proposed year is recorded as \(2017\); no individual proposer is identified.
\item Supplied source citation: \emph{O-Minimality and its Applications to Number Theory and Analysis} (2018), Oberwolfach Reports, pp.~1349--1420, DOI \texttt{10.4171/owr/2017/22}. \url{https://doi.org/10.4171/owr/2017/22}. The background describes these reports as documenting mathematical workshops and their open-problem discussions.
\item Supplied literature source, labeled ``Oberwolfach Reports 2018'': \url{https://doi.org/10.4171/owr/2017/22}.
\item Supplied literature source, attributed to Rainer and labeled ``arc-smooth functions on closed sets'': \url{https://arxiv.org/abs/1801.08335}.
\item Literature assessment dated exactly \texttt{2026-08-21}: the supplied curation labels the problem partially solved and reports equivalence of the smooth and ultradifferentiable classes under uniform cusp or cone hypotheses and in several subanalytic cases, while describing the general compact-set question as open. This is supplied dated curation, not an independently verified assessment. Its hypotheses and cited claims require independent verification.
\end{itemize}
\end{document}
